\section{Complete fixed coefficients of the ten-variable equation}
\label{app:compact-coefficients}

This appendix defines the four integers $\beta,K_0,M_C,M_F$ used in the
degree-190 equation.  Every index, digit, tile, finite set, sum, product and
floor appearing below is part of a fixed finite recipe.  None introduces
an additional quantified unknown.  The recipe does not inspect whether
ZFC is consistent and does not search for a polynomial zero.

\paragraph{Provenance and notation.}
These are the unchanged BFS943/radix67 constants retained by source~172
at the inspected ProveIt commit
\texttt{d123610bc98fef2006bff1c9466bc3f1221003e7}.
The archived finite specification is
\path{compact/source/upstream172/inherited/strict/upstream_lucas/parent/proof_sources/indexed_constants.tex}.
The operative reconstruction is
\path{compact/source/rebuild_coefficients.py}.
Only the fixed-coefficient recipe is reproduced here; the historical
equation surrounding that recipe is not part of the present definition.
Its constant $K$ is denoted $K_0$ to match the body of this paper.
All other notation in this appendix is local to the finite recipe.

\subsection{The four integers}

Put
\[
 \rho=2^{67},\qquad B=\rho^{8807},\qquad
 \beta=B-1=2^{590069}-1,\qquad
 G_m(x)=\sum_{j=0}^{m-1}x^j.
\]
Define the finite exponent set
\begin{align*}
 \mathcal E={}&\{0,\ldots,942\}\cup\{983,\ldots,992\}
                 \cup\{1003,\ldots,1011\}\\
 &\cup\{1013,\ldots,1031\}
                 \cup\{1063,1064,1095,2195,3296\}.
\end{align*}
Then the other three fixed integers are
\begin{align*}
 M_C={}&B-1-\sum_{e\in\mathcal E\setminus\{1\}}\rho^e
                     -2\rho^{1063},\\
 M_F={}&B-1+4+(2^{51}-2)(\rho^{4393}+\rho^{5495})\\
       &+\rho^{1013}G_{10}(\rho)+\rho^{1065}G_9(\rho)
                     +\rho^{2195}+\rho^{3296}\\
       &+(2^{65}-1)\rho^{8792}G_{13}(\rho)
                     +(2^{19}-1)\rho^{8805},\\
 K_0={}&\rho^{30}+\rho^{62}+\rho^{1100}+\rho^{2201}
                     +\rho^{8849}\\
       &+(\rho^{4393}+\rho^{5495})\Phi(1/\rho),
\end{align*}
where
\begin{align*}
 \Phi(x)={}&(1+2^{51})G_{943}(x)
           -2^{10}\sum_{i=0}^{942}v(t_i)x^i
           +\sum_{(b,e,m)\in\mathcal C}2^b x^eG_m(2x),\\
 \mathcal C={}&\{(10,1023,8),(18,1013,9),(27,1003,8),
                  (35,983,9),\\
             &\quad(44,1031,1),(45,1022,1),(46,1011,1),
                  (47,992,1),\\
             &\quad(48,1095,1),(49,2195,1),(50,3296,1)\}.
\end{align*}
The tiles $t_i$ and their integer values $v(t_i)$ are specified next.
The largest exponent in $\Phi$ is 3296, smaller than both 4393 and 5495.
Consequently the displayed rational intermediate in $K_0$ expands into
integer powers of $\rho$.  The four resulting integers are positive.

\subsection{The 943 indexed fixed tiles}

Use the 430 actions $(u_j,v_j,d_j)$ defined in the next subsection.  Set
\begin{align*}
 h(q,b)&=2(q+1)+b,\\
 \mathcal R&=\{u_j:d_j=1,\ 0\leq j<430\},
 &\mathcal L&=\{u_j:d_j=0,\ 0\leq j<430\},\\
 r_q&=6+|\mathcal R\cap\{0,\ldots,q-1\}|,
 &l_q&=205+|\mathcal L\cap\{0,\ldots,q-1\}|.
\end{align*}
An interval ending at $-1$ is empty.  The following table defines every
$t_i=(w,n,e,s)$ for $0\leq i\leq942$.  The index ranges and the actions
are fixed data, not parameters of the Diophantine equation.

\begingroup
\renewcommand{\arraystretch}{1.35}
\begin{longtable}{@{}p{0.45\linewidth}p{0.51\linewidth}@{}}
\toprule
Index $i$ and fixed range & Tile $t_i=(w,n,e,s)$\\
\midrule
\endfirsthead
\toprule
Index $i$ and fixed range & Tile $t_i=(w,n,e,s)$\\
\midrule
\endhead
$0$ & $(2,435,3,1)$\\
$1$ & $(5,0,0,0)$\\
$j+2$, where $j=2q+b<430$ &
 $((1-d_j)l_{u_j},\ h(q,b),\ d_jr_{u_j},\ v_j)$\\
$432+b$, $b\in\{0,1\}$ & $(0,h(215,b),0,h(215,b))$\\
$434+b$, $b\in\{0,1\}$ & $(0,b,0,b)$\\
$424+2r_q+b$, $q\in\mathcal R$, $b\in\{0,1\}$ &
 $(r_q,b,0,h(q,b))$\\
$424+2l_q+b$, $q\in\mathcal L$, $b\in\{0,1\}$ &
 $(0,b,l_q,h(q,b))$\\
$924+j$, $0\leq j<4$ &
 $(0,j+430\lfloor j/2\rfloor,0,438)$\\
$928$ & $(2,434,2,0)$\\
$929$ & $(3,436,4,4)$\\
$930$ & $(4,436,4,0)$\\
$497+t$, $t\in\{434,435,436\}$ & $(1,438,1,t)$\\
$934+3a+b-6\lfloor a/4\rfloor$,\newline
 $a\in\{0,1,4\}$, $b\in\{0,1,2\}$ & $(a,437,b,437)$\\
\bottomrule
\end{longtable}
\endgroup

For each tile $t=(w,n,e,s)$, define
\[
 v(t)=w+2^8n+2^{17}e+2^{25}s+127\cdot2^{34}.
\]

\subsection{A closed decoder for the 430 actions}

For fixed nonnegative integers define
\[
 D_b(A,k)=\left\lfloor\frac{A}{b^k}\right\rfloor\bmod b,
 \qquad
 a\bmod b=a-b\left\lfloor\frac{a}{b}\right\rfloor.
\]
Let $\mathcal A$ be the single fixed integer specified in the next
subsection.  First split it into three fixed blocks:
\[
 N=\mathcal A\bmod2^{636},\qquad
 T_0=\left\lfloor\frac{\mathcal A}{2^{636}}\right\rfloor
          \bmod2^{1675},\qquad
 E_0=\left\lfloor\frac{\mathcal A}{2^{2311}}\right\rfloor.
\]
For $0\leq q<215$, $b\in\{0,1\}$ and $j=2q+b$, put
\begin{align*}
 \nu_j&=D_2(N,j),
 &\sigma_j&=\sum_{k=0}^{j-1}\nu_k,\\
 m_q&=D_2(N,430+q),
 &\tau_q&=\sum_{k=0}^{q-1}m_k,\\
 F_q&=14(1-m_q)+m_qD_{16}(E_0,\tau_q),\\
 u_j&=\nu_j(1+\sigma_j)
              +(1-\nu_j)D_{216}(T_0,j-\sigma_j),\\
 v_j&=D_2(F_q,2b),
 &d_j&=D_2(F_q,2b+1).
\end{align*}
Empty sums are zero.  In particular,
$(v_0,d_0,u_0)=(1,1,1)$; the constructor tile uses this actual first
target, including its state renaming.

\subsection{The 436-character numeral}

The alphabet, in increasing order of digit values $0,\ldots,61$, is
\begin{center}
\texttt{0123456789ABCDEFGHIJKLMNOPQRSTUVWXYZabcdefghijklmnopqrstuvwxyz}.
\end{center}
Concatenate the following eight lines without spaces or line breaks.
Interpret their 436 characters as digit values $e_0,\ldots,e_{435}$ in
that alphabet, and set
\[
 \mathcal A=\sum_{j=0}^{435}e_j62^{435-j}.
\]
\begin{lstlisting}[basicstyle=\small\ttfamily,breaklines=false]
BxEnhvy6NipzkWmROo0GsA3uK3twWLxscLq20UTumy5Cdu461LwD6g7RDNzF8Y
iYDFTpB8ehAypd7d6OsOL4APamnJeNXQypUezCmW9n4oXP6XfgmiWuAWZ8ZJ7U
OP7gu4kVdlGZqZHULQMar1dXNeJW4i7oNVI8usp8BlovjjZbwPCovZI328fRUq
4q92R81s5GjWQaQfGhkq9UYMK3rRMaBmsOIBkJqc2cj2VRIZ6mjzkBIOkW9SjZ
FzeA6J5u4j3sBL3SdgwXzM9mHvQivu8EJ1TBTpRhHDqZIbmYgwQ8VqTAFotl6F
We1YC554GKrVHsA8oUGKSA6F6k5KREBSD0Tgd7bFALY84LLbpw574BQyvdoHrl
0Oe35uwDf15DpT6VJWZBjSbHPpPAINYTte0eUxnTFTKBlGj8Am6sb1gLldh8WZ
8P
\end{lstlisting}

These data close the coefficient definitions: they determine the actions,
then all 943 tiles, then $\Phi$, and finally the four integers.
The supplied reconstruction checks the exact integers against the
authoritative payloads, independently of the gate evaluator.  The finite
decoder, tile indices, exponent ranges and numeral above agree with that
reconstruction.  This appendix supplies fixed coefficients; the proof
of their connection to refutations is the separately identified inherited
calibration.
